\documentclass[11pt]{amsart}

\usepackage[T1]{fontenc}
\usepackage{lmodern}
\usepackage{microtype}
\usepackage{amsmath,amssymb,amsthm}
\usepackage{booktabs}
\usepackage[hidelinks]{hyperref}

\hypersetup{
  pdftitle={A Negative Answer to Problem 8.8 of Galetto, Montano, and Wellner}
}

\newtheorem{theorem}{Theorem}
\newtheorem{proposition}[theorem]{Proposition}
\newtheorem{lemma}[theorem]{Lemma}

\newcommand{\F}{\mathbb F}
\newcommand{\Q}{\mathbb Q}
\newcommand{\VR}{\operatorname{VR}}
\newcommand{\Aut}{\operatorname{Aut}}
\newcommand{\Ind}{\operatorname{Ind}}
\newcommand{\diameter}{\operatorname{diam}}
\newcommand{\stype}[1]{\texttt{#1}}

\title[A negative answer to Problem 8.8]
{A Negative Answer to Problem~8.8 of Galetto, Monta\~no, and Wellner}

\date{}

\subjclass[2020]{Primary 05E18; Secondary 05E45, 55N31}
\keywords{Vietoris--Rips complex, hypercube, hyperoctahedral group,
multiplicity-free representation, Lefschetz number}

\begin{document}

\begin{abstract}
Let $K=\VR(Q_6;4)$ and $G=\Aut(Q_6)$. We prove that, over every field
$\Bbbk$ of characteristic zero, and conditional on the Betti numbers of $K$
computed in \cite[Appendix~C.3]{GMW},
\[
\dim_{\Bbbk}H_7(K;\Bbbk)^G+
\dim_{\Bbbk}H_{15}(K;\Bbbk)^G=3.
\]
Thus the trivial representation occurs with multiplicity at least two in one
of these two homology groups. This gives a negative answer to Problem~8.8 of
Galetto, Monta\~no, and Wellner. The proof is an exact Lefschetz computation
on the $65$ conjugacy classes of the hyperoctahedral group.
\end{abstract}

\maketitle
\raggedbottom

\section{The counterexample}

Let $Q_n$ be the $n$-dimensional hypercube graph, and identify it throughout
with its vertex set $\{0,1\}^n$ carrying the Hamming metric $d$, which is the
graph metric of $Q_n$. In particular $\Aut(Q_n)$ is simultaneously the
automorphism group of the graph and the isometry group of the metric space,
and the Vietoris--Rips complex below is taken on that vertex set:
\[
X^{n,r}=\VR(Q_n;r)
=\{\sigma\subseteq Q_n:d(x,y)\le r\text{ for all }x,y\in\sigma\}.
\]
The hyperoctahedral group
\[
\mathfrak H_n=\Aut(Q_n)\cong \F_2^n\rtimes\mathfrak S_n
\]
acts simplicially on $X^{n,r}$. Galetto, Monta\~no, and Wellner proved
multiplicity-freeness of the nonzero homology representations at scales
$r\le3$ and $r=n-1$, and asked whether $H_i(X^{n,r})$ is multiplicity-free
for all $n,r,i$ \cite[Problem~8.8]{GMW}. The smallest pair with
$4\le r\le n-2$ is $(n,r)=(6,4)$. They also computed
\begin{equation}\label{eq:homology}
\widetilde H_i(X^{6,4};\Q)\cong
\begin{cases}
\Q^{239},&i=7,\\
\Q^{14},&i=15,\\
0,&i\notin\{7,15\}.
\end{cases}
\end{equation}
See \cite[Appendix~C.3]{GMW}. Equation~\eqref{eq:homology} is a machine
computation of theirs which we do not reproduce; it is the only computational
input taken as given here, and Theorem~\ref{thm:main} is conditional on it,
both for the vanishing of $\widetilde H_i(X^{6,4};\Q)$ outside $\{7,15\}$ and
for the parity of the two surviving degrees.

\begin{theorem}\label{thm:main}
Let $K=X^{6,4}$ and $G=\mathfrak H_6$. For every field $\Bbbk$ of
characteristic zero,
\[
\dim_{\Bbbk}H_7(K;\Bbbk)^G+
\dim_{\Bbbk}H_{15}(K;\Bbbk)^G=3.
\]
Consequently, at least one of the $G$-representations $H_7(K;\Bbbk)$ and
$H_{15}(K;\Bbbk)$ is not multiplicity-free.
\end{theorem}

The argument identifies the counterexample at $(n,r)=(6,4)$ but does not
determine which of the two degrees contains the repeated trivial constituent.

\section{A fixed-orbit Lefschetz formula}

For $g\in G$, define
\[
\Psi(g)=\sum_{i\ge0}(-1)^{i+1}
\operatorname{tr}\!\left(g\mid\widetilde H_i(K;\Q)\right).
\]
Let $\mathcal O_g$ be the set of $g$-orbits on $Q_6$ with internal Hamming
diameter at most $4$. Define a graph $F_g$ on $\mathcal O_g$ by joining
$O,O'$ when
\[
\diameter(O\cup O')>4.
\]
For a graph $F$, write
\[
I(F;t)=\sum_{S\in\Ind(F)}t^{|S|}
\]
for its independence polynomial.

\begin{lemma}\label{lem:fixed-orbit}
For every $g\in G$,
\[
\Psi(g)=I(F_g;-1).
\]
\end{lemma}

\begin{proof}
A simplex fixed setwise by $g$ is a union of vertex orbits. Such a union is a
simplex of $K$ exactly when the selected orbits form an independent set of
$F_g$.

Suppose a fixed simplex $\sigma$ has $s$ vertices and is the union of $c$
vertex orbits. The induced permutation of its vertices has $c$ cycles, hence
orientation sign $(-1)^{s-c}$. Since $\sigma$ lies in chain degree $s-1$, its
contribution to the reduced Lefschetz number is
\[
(-1)^{s-1}(-1)^{s-c}=(-1)^{c+1}.
\]
The empty simplex in degree $-1$ contributes $-1$, which is the same formula
for $c=0$. The Hopf trace formula, applied to the augmented chain complex
(indexed by $i\ge-1$) and its reduced homology, therefore gives
\[
\sum_{i\ge-1}(-1)^i\operatorname{tr}
 \left(g\mid\widetilde H_i(K;\Q)\right)
=-\sum_{S\in\Ind(F_g)}(-1)^{|S|}.
\]
Here $\widetilde H_{-1}(K)=0$ because $K$ is nonempty, so the left-hand side
reduces to the sum over $i\ge0$, which equals $-\Psi(g)$.
Multiplying by $-1$ proves the claim.
\end{proof}

\section{Exact class computation}

Write $g=(a,\pi)\in\F_2^6\rtimes\mathfrak S_6$, acting by
$x\mapsto a+\pi x$. For each cycle $C$ of $\pi$, the parity
$\sum_{j\in C}a_j\in\F_2$ is invariant under conjugacy. These parities,
together with the cycle lengths, form a complete invariant: conjugacy classes
are indexed by bipartitions $(\alpha;\beta)$ of $6$, where $\alpha$ records
the even-parity cycles and $\beta$ the odd-parity cycles
\cite[Section~5.2]{GMW}.

If $m_j(\lambda)$ is the multiplicity of $j$ in a partition $\lambda$, set
\[
z_\lambda=\prod_{j\ge1}j^{m_j(\lambda)}m_j(\lambda)!.
\]
A signed cycle of length $j$ contributes a factor $2j$ to the centralizer,
and signed cycles of the same length and parity may be permuted. Hence
\begin{equation}\label{eq:class-size}
|C_{\alpha,\beta}|=
\frac{2^6 6!}{2^{\ell(\alpha)+\ell(\beta)}z_\alpha z_\beta}.
\end{equation}
There are $65$ bipartitions of $6$.

For each type $(\alpha;\beta)$, choose disjoint cycles on consecutive
coordinate blocks, take $a=0$ on every $\alpha$-cycle, and put a single $1$
in $a$ on every $\beta$-cycle. Lemma~\ref{lem:fixed-orbit} then reduces the
class value to an exact graph computation. If $J(F)=I(F;-1)$ and $N[v]$ is
the closed neighborhood of $v$, then
\begin{equation}\label{eq:recurrence}
J(F)=J(F-v)-J(F-N[v]),\qquad J(\varnothing)=1,
\end{equation}
and $J$ is multiplicative over connected components.

Appendix~\ref{app:table} lists the resulting value for every class. Together
with \eqref{eq:class-size} and \eqref{eq:recurrence} the table is a
self-contained certificate: the recipe of the preceding paragraph fixes a
representative of each of the $65$ classes, the $g$-orbits on $Q_6$ and the
graph $F_g$ are read off from that representative, \eqref{eq:recurrence}
evaluates $J(F_g)$ in integer arithmetic, and \eqref{eq:class-size} supplies
the class size. Both sums below are then finite sums of $65$ integers, which a
reader can recompute directly.

The identity class also checks the one input we take as given. Since $7$ and
$15$ are odd, \eqref{eq:homology} forces $\Psi(1)=239+14=253$, and the value
computed here from $F_1$ by \eqref{eq:recurrence} is the entry $253$ at type
\stype{111111|-} in Appendix~\ref{app:table}. Thus the reduced Euler
characteristic $\widetilde\chi(X^{6,4})=-253$ is recovered from the
fixed-orbit data alone, and an error in the Betti numbers of
\cite[Appendix~C.3]{GMW} would have to leave their alternating sum intact.

\begin{proposition}\label{prop:average}
The class function $\Psi$ satisfies
\[
\sum_{g\in G}\Psi(g)=138240,
\qquad
\frac1{|G|}\sum_{g\in G}\Psi(g)=3.
\]
\end{proposition}

\begin{proof}
Using \eqref{eq:class-size} and the values in Appendix~\ref{app:table}, exact
summation gives
\[
\sum_{\alpha,\beta}|C_{\alpha,\beta}|=46080=2^6 6!,
\qquad
\sum_{\alpha,\beta}|C_{\alpha,\beta}|\Psi(\alpha;\beta)=138240.
\]
Division gives the stated average.
\end{proof}

\begin{proof}[Proof of Theorem~\ref{thm:main}]
For a finite-dimensional $G$-representation $V$ over $\Q$, character
averaging gives
\[
\frac1{|G|}\sum_{g\in G}\operatorname{tr}(g\mid V)=\dim_{\Q}V^G.
\]
By \eqref{eq:homology}, the only nonzero reduced homology groups occur in the
odd degrees $7$ and $15$. Moreover $H_i(K;\Bbbk)=\widetilde H_i(K;\Bbbk)$ for
$i\ge1$, so the two degrees in question are unaffected by the passage between
reduced and unreduced homology. Proposition~\ref{prop:average} therefore
yields
\[
3=\dim_{\Q}H_7(K;\Q)^G+
  \dim_{\Q}H_{15}(K;\Q)^G.
\]
At least one summand is at least $2$. Since $\Q[G]$ is semisimple, this
summand is the multiplicity of the trivial representation.

For any characteristic-zero field $\Bbbk$, extension of scalars gives
\[
H_i(K;\Bbbk)\cong H_i(K;\Q)\otimes_{\Q}\Bbbk.
\]
Invariants commute with this extension because they are the image of the
Reynolds idempotent $|G|^{-1}\sum_{g\in G}g$. Thus the same equality and
conclusion hold over $\Bbbk$.
\end{proof}

\appendix
\section{Conjugacy-class certificate}\label{app:table}

Partitions are written by concatenating their parts, and \stype{-} denotes
the empty partition. Thus \stype{22|11} denotes $((2,2);(1,1))$; the first
partition records even-parity cycles and the second odd-parity cycles.

\begin{table}[htbp]
\centering
\scriptsize
\setlength{\tabcolsep}{3.2pt}
\begin{tabular}{@{}lrr@{\qquad}lrr@{\qquad}lrr@{}}
\toprule
Type & $|C|$ & $\Psi$ & Type & $|C|$ & $\Psi$ & Type & $|C|$ & $\Psi$\\
\midrule
\stype{-|111111} & 1 & 1 & \stype{-|21111} & 30 & 1
  & \stype{-|2211} & 180 & 17 \\
\stype{-|222} & 120 & 1 & \stype{-|3111} & 160 & 1
  & \stype{-|321} & 960 & 1 \\
\stype{-|33} & 640 & 1 & \stype{-|411} & 720 & 1
  & \stype{-|42} & 1440 & 1 \\
\stype{-|51} & 2304 & 1 & \stype{-|6} & 3840 & 1
  & \stype{1|11111} & 6 & 1 \\
\stype{1|2111} & 120 & 5 & \stype{1|221} & 360 & 17
  & \stype{1|311} & 480 & 1 \\
\stype{1|32} & 960 & 5 & \stype{1|41} & 1440 & 5
  & \stype{1|5} & 2304 & 1 \\
\stype{11|1111} & 15 & 53 & \stype{11|211} & 180 & 5
  & \stype{11|22} & 180 & 5 \\
\stype{11|31} & 480 & 5 & \stype{11|4} & 720 & 5
  & \stype{111|111} & 20 & 53 \\
\stype{111|21} & 120 & 9 & \stype{111|3} & 160 & 5
  & \stype{1111|11} & 15 & 9 \\
\stype{1111|2} & 30 & 9 & \stype{11111|1} & 6 & 9
  & \stype{111111|-} & 1 & 253 \\
\stype{2|1111} & 30 & 3 & \stype{2|211} & 360 & -1
  & \stype{2|22} & 360 & 3 \\
\stype{2|31} & 960 & 3 & \stype{2|4} & 1440 & -1
  & \stype{21|111} & 120 & 3 \\
\stype{21|21} & 720 & 3 & \stype{21|3} & 960 & 3
  & \stype{211|11} & 180 & 7 \\
\stype{211|2} & 360 & 3 & \stype{2111|1} & 120 & 7
  & \stype{21111|-} & 30 & 11 \\
\stype{22|11} & 180 & 9 & \stype{22|2} & 360 & 1
  & \stype{221|1} & 360 & 9 \\
\stype{2211|-} & 180 & 29 & \stype{222|-} & 120 & 3
  & \stype{3|111} & 160 & -1 \\
\stype{3|21} & 960 & 3 & \stype{3|3} & 1280 & 5
  & \stype{31|11} & 480 & 3 \\
\stype{31|2} & 960 & 3 & \stype{311|1} & 480 & 3
  & \stype{3111|-} & 160 & 7 \\
\stype{32|1} & 960 & 1 & \stype{321|-} & 960 & 5
  & \stype{33|-} & 640 & 13 \\
\stype{4|11} & 720 & 3 & \stype{4|2} & 1440 & 3
  & \stype{41|1} & 1440 & 3 \\
\stype{411|-} & 720 & 7 & \stype{42|-} & 1440 & 5
  & \stype{5|1} & 2304 & -1 \\
\stype{51|-} & 2304 & 3 & \stype{6|-} & 3840 & 3
  & & & \\
\bottomrule
\end{tabular}
\caption{Exact values of $\Psi(g)=I(F_g;-1)$ on all conjugacy classes.}
\label{tab:certificate}
\end{table}

\begin{thebibliography}{9}

\bibitem{GMW}
F.~Galetto, J.~Monta\~no, and Z.~Wellner,
\emph{Homology of Vietoris--Rips complexes of hypercube graphs via group
actions},
\href{https://arxiv.org/abs/2606.20784}{arXiv:2606.20784v1}, 2026.

\end{thebibliography}

\end{document}